% Frozen isolated candidate proof.  This file is not a live-book input.
% Production Rethlas verification is pending.

We divide the proof into source-legality, a countable-factor comparison
sheaf, and a fixed-neighbourhood regularization.

First consider \(Z\) alone.  Its topology satisfies
\cite[\S3.1.1]{CLD25}: local compactness and the clopen-basis hypothesis
give a basis of relatively compact open subsets, hence a basis of open
subsets that are countable at infinity.  Locally closed subsets satisfy the
domain axioms of \cite[\S3.1.2]{CLD25}.  Indeed, the first three axioms are
immediate from the definition, and every locally closed subset is itself a
locally compact space with a clopen basis; its compact clopen
neighbourhoods form the required \(G\)-cover by compact domains.

The locally constant real-valued functions form a \(G\)-sheaf.  To check
the only nonformal point, suppose compatible locally constant functions are
glued across a \(G\)-cover.  At a given point, use the finite subfamily in
the definition of a \(G\)-cover.  Intersect ordinary neighbourhoods on
which the finitely many restrictions are constant.  Their union contains
a neighbourhood of the point in the domain, and compatibility makes the
glued function constant there.

A locally constant map from a compact domain has finite image.  In the
standing \((\mathbf Q,\mathbf R)\)-convention of
\cite[\S1.8.2]{CLD25}, every finite subset of \(\mathbf R^m\) is a compact
polyhedron.  Parlinear composition preserves local constancy.  The inverse
image of any subset under a locally constant map is clopen in its domain,
so in particular the inverse image of a parlinear subspace is a domain.
This verifies all the axioms in \cite[\S3.1.3]{CLD25}.

Take the affine presheaf \(L_Z\) equal to this parlinear sheaf.  It contains
the constants, ordinary local gluing preserves local constancy, and a
function is parlinear if and only if it is \(G\)-locally affine.  Thus the
two affinage axioms of \cite[\S6.1.1]{CLD25} hold.  Consequently the data
in the statement make \(Z\) an affine \(G\)-tropical space of tropical
dimension zero.

The product with the standard affine space \(\Omega\) is legitimate by
\cite[\S\S9.1.1--9.1.4]{CLD25}.  Product affine functions are locally sums
of locally constant functions of \(z\) and affine functions of \(x\).
It follows from \cite[Definition~6.2.1]{CLD25} that, after a clopen
shrinking in the \(Z\)-factor, a smooth function is a smooth function of
the Euclidean variable alone, and conversely.  The formula for
\(\mathrm d'\mathrm d''\) and positivity are therefore the ordinary
Hessian formula and positive-semidefiniteness on the Euclidean factor.
This proves source-legality and the asserted description of the smooth psh
generators.

Define \(\mathscr Q\) to be the sheaf of functions with values in
\(\mathbf R\cup\{-\infty\}\) which locally have the following form.  There
are a compact clopen \(C\subseteq Z\), an open ball
\(B\Subset\Omega\), a compact metrizable space \(Y\) with a clopen basis, a
continuous surjection \(q:C\to Y\), and an upper semicontinuous function
\(\bar v:Y\times B\to\mathbf R\cup\{-\infty\}\) such that
\begin{equation}
\label{eq:p071-stone-local-factor}
                  v(z,x)=\bar v(q(z),x),
\end{equation}
and, for every \(y\in Y\), \(x\mapsto\bar v(y,x)\) is either finite-valued
convex or identically \(-\infty\).  This is a sheaf because the condition
is local.  Every smooth psh function belongs to
\(\mathscr Q\): shrink until it is independent of \(z\), take \(Y\) to be
a point, and use positivity of its Euclidean Hessian.

We next prove that \(\mathscr Q\) is stable under decreasing sequential
limits.  Let \(v_n\downarrow v\) on an open subset \(W\) of \(X\), with
every \(v_n\in\mathscr Q(W)\), and fix a point.  Choose a compact clopen
\(C\subseteq Z\) and a ball \(B_0\Subset\Omega\) such that
\begin{equation}
\label{eq:p071-stone-compact-core}
                         C\times\overline{B_0}\subseteq W.
\end{equation}
For each
fixed \(n\), compactness gives a finite cover of
\(C\times\overline{B_0}\) by product sets
\begin{equation}
\label{eq:p071-stone-factor-cover}
                       C_{n,a}\times B_{n,a},
\end{equation}
where \(C_{n,a}\) is compact clopen in \(C\), \(B_{n,a}\) is open, and on this
product \(v_n\) factors through a compact metrizable zero-dimensional
quotient \(q_{n,a}:C_{n,a}\to Y_{n,a}\).

Choose a countable clopen base \(\mathcal D_{n,a}\) of every \(Y_{n,a}\).
Let \(\mathcal B\) be the Boolean subalgebra of the clopen subsets of \(C\)
generated by all the sets \(C_{n,a}\) and all the sets
\begin{equation}
\label{eq:p071-stone-boolean-generators}
                  q_{n,a}^{-1}(D),
                  \qquad D\in\mathcal D_{n,a}.
\end{equation}
viewed as clopen subsets of \(C\).  There are countably many generators.
Declare two points of \(C\) equivalent when they belong to the same members
of \(\mathcal B\), and let
\begin{equation}
\label{eq:p071-stone-common-quotient}
                  q:C\longrightarrow
                  Y:=q(C)\subseteq\{0,1\}^{\mathcal B}.
\end{equation}
Every coordinate of \(q\) is the characteristic function of a clopen set,
so \(q\) is continuous and is surjective by the definition of \(Y\).
The space \(Y\) is compact, metrizable, and zero-dimensional.

Every \(v_n|_{C\times B_0}\) is constant on the fibres of
\(q\times\operatorname{id}_{B_0}\).  Indeed, if \(q(z)=q(z')\) and
\(x\in B_0\), choose a member of the finite cover
\eqref{eq:p071-stone-factor-cover} containing
\((z,x)\).  Membership in \(C_{n,a}\) is recorded in \(\mathcal B\), so
\(z'\in C_{n,a}\).  Membership in the inverse images of the countable
clopen base separates the points of \(Y_{n,a}\), hence
\(q_{n,a}(z)=q_{n,a}(z')\), and the local factorization gives
\(v_n(z,x)=v_n(z',x)\).  It follows that the decreasing limit \(v\) factors
as
\begin{equation}
\label{eq:p071-stone-limit-factor}
                         v(z,x)=\bar v(q(z),x).
\end{equation}

The function \(\bar v\) is upper semicontinuous.  Here is the quotient
argument, included to avoid an unsupported product-quotient assertion.
For \(t\in\mathbf R\), the set \(\{v<t\}\) is open because
\begin{equation}
\label{eq:p071-stone-sublevel-union}
                       \{v<t\}=\bigcup_n\{v_n<t\}.
\end{equation}
If \(\bar v(y,x)<t\), then, for every \(z\in q^{-1}(y)\), choose open
sets \(N_z\ni z\) and \(D_z\ni x\) whose product lies in \(\{v<t\}\).
A finite subcover of the compact fibre gives
\(N=\bigcup_zN_z\) and \(D=\bigcap_zD_z\), where the union and intersection
are now over that finite subcover; then
\(N\times D\subseteq\{v<t\}\).  The map \(q\) is closed because its source
is compact and its target is Hausdorff.  Therefore
\begin{equation}
\label{eq:p071-stone-quotient-neighbourhood}
                       Y\setminus q(C\setminus N)
\end{equation}
is an open neighbourhood of \(y\), and its product with \(D\) is contained
in \(\{\bar v<t\}\).  Thus all strict sublevel sets of \(\bar v\) are open.

We include the extended-value point in the fibrewise convexity argument.
A locally extended-real convex function
\(h:B_0\to\mathbf R\cup\{-\infty\}\) is either finite everywhere or
identically \(-\infty\).  On a sufficiently small ball of local convexity,
if \(h(a)=-\infty\), every point of a smaller ball about \(a\) lies
strictly between \(a\) and another point of the first ball, so convexity
forces the smaller ball to have value \(-\infty\).  If \(h(a)\) is finite,
the same midpoint argument excludes \(-\infty\) on a smaller ball.  Hence
the \(-\infty\)-locus and its complement are open; connectedness of \(B_0\)
proves the dichotomy.  In the finite case, subdivision of line segments
shows that local convexity implies convexity on \(B_0\).

Apply this observation to \(x\mapsto v_n(z,x)\), for a fixed \(z\).  If one
term is identically \(-\infty\), all later terms and the limit are as well.
Otherwise all terms are finite and convex.  Passing to the limit in their
convexity inequalities shows that the limit is extended-real convex.  If
it is \(-\infty\) at one point, then, for any other point, use a slightly
longer segment inside the open convex ball and the convexity inequality to
force the value at that point to be \(-\infty\).  This is the elementary
argument used in \cite[Proposition~8.3.2]{CLD25}.  Otherwise the limit is
finite and convex.  By \eqref{eq:p071-stone-limit-factor}, this is exactly
the required statement for every fibre of \(\bar v\).  We have proved that
\(\mathscr Q\) is stable under decreasing limits.

The sheaf \(\mathscr Q\) contains every smooth psh generator and is stable
under the defining operation of \(\operatorname{Psh}_X\).  Minimality in
\cite[Definition~8.1.1]{CLD25} therefore gives
\begin{equation}
\label{eq:p071-stone-minimality}
                         \operatorname{Psh}_X\subseteq\mathscr Q.
\end{equation}

It remains to prove the reverse inclusion into the first stage, and this is
where the representatives are re-selected on one fixed neighbourhood.  Let
\(v\in\mathscr Q\) have a representation
\eqref{eq:p071-stone-local-factor} on \(C\times B_0\).  Suppose first that
\(d>0\), fix any ball \(B\) with \(\overline B\subseteq B_0\), and choose a
ball \(B_1\) such that
\begin{equation}
\label{eq:p071-stone-nested-balls}
                 \overline B\subseteq B_1,
                 \qquad \overline{B_1}\subseteq B_0
\end{equation}
and positive radii \(r_k\downarrow0\) such that
\(\overline{B_1}+r_1\overline{\mathbb B}\subseteq B_0\).  Give \(Y\) a
compatible metric.  Compactness and the clopen-basis property give nested
finite clopen partitions \(\mathcal P_k\) whose mesh tends to zero: at
stage \(k\), take a finite clopen cover of mesh less than \(2^{-k}\) and
refine it together with all preceding covers.  For
\(P\in\mathcal P_k\) and \(x\in B_1\), put
\begin{equation}
\label{eq:p071-stone-upper-envelope}
 a_{k,P}(x)=
 \max\left\{-k,
   \sup_{\substack{y\in P\\ |\xi|\leq r_k}}
          \bar v(y,x+\xi)\right\}.
\end{equation}
Upper semicontinuity bounds \(\bar v\) from above on the compact set
\(Y\times(\overline{B_1}+r_1\overline{\mathbb B})\).  Thus the supremum in
\eqref{eq:p071-stone-upper-envelope} is finite from above, and the
truncation makes \(a_{k,P}\) real-valued.  It is convex, being the maximum
of a constant and a supremum of translates of convex functions, and hence
is continuous on \(B_1\).  If \(P_k(y)\) denotes the member of
\(\mathcal P_k\) containing
\(y\), define
\begin{equation}
\label{eq:p071-stone-piecewise-envelope}
                         a_k(y,x)=a_{k,P_k(y)}(x).
\end{equation}
The refinement of the partitions, the decrease of \(r_k\), and the lower
truncations \(-k\) give
\begin{equation}
\label{eq:p071-stone-envelope-monotonicity}
                          a_{k+1}\leq a_k.
\end{equation}
Moreover,
\begin{equation}
\label{eq:p071-stone-envelope-convergence}
                          a_k\downarrow\bar v
\end{equation}
pointwise on \(Y\times B_1\).  To check the only nonformal assertion, fix
\((y,x)\).  The lower bound follows by taking \(\xi=0\) and the same \(y\)
in \eqref{eq:p071-stone-upper-envelope}.  If \(\bar v(y,x)\) is finite,
upper semicontinuity makes
\(\bar v\) smaller than \(\bar v(y,x)+\varepsilon\) on a product
neighbourhood; eventually \(P_k(y)\) lies in its \(Y\)-factor and
\(r_k\) lies in its Euclidean factor.  If \(\bar v(y,x)=-\infty\), apply
the same argument with an arbitrary finite upper level, also taking \(k\)
large enough that \(-k\) lies below that level, and then let the level tend
to \(-\infty\).  This proves
\eqref{eq:p071-stone-envelope-convergence}, including the polar case.

Choose positive numbers \(\eta_k\) with
\(\sum_k\eta_k<\infty\).  Let \(\rho\) be a nonnegative smooth
probability kernel on \(\mathbf R^d\), supported in the unit ball and of
barycentre zero.  Since every \(a_{k,P}\) is a finite continuous convex
function and \(\mathcal P_k\) is finite, one can choose
\(\varepsilon_k>0\), with the convolution defined on \(B\), such that
\begin{equation}
\label{eq:p071-stone-mollification-error}
  0\leq
  (a_{k,P}*\rho_{\varepsilon_k})(x)-a_{k,P}(x)
  \leq\eta_k
  \quad(P\in\mathcal P_k,\ x\in B).
\end{equation}
The lower inequality is Jensen's inequality; the upper inequality is the
uniform convergence of mollifications on the compact set
\(\overline B\), applied to the finite family indexed by
\(\mathcal P_k\).

On \(q^{-1}(P)\times B\), define
\begin{equation}
\label{eq:p071-stone-smoothed-envelope}
                 w_k(z,x)=
                 (a_{k,P}*\rho_{\varepsilon_k})(x).
\end{equation}
The sets \(q^{-1}(P)\) form a finite clopen partition of \(C\), and each is
clopen in \(Z\).  Hence
\(w_k\) is smooth for the stated product affine structure.  Each branch in
\eqref{eq:p071-stone-smoothed-envelope} is smooth convex, so
\(\mathrm d'\mathrm d''w_k\geq0\).  Equations
\eqref{eq:p071-stone-envelope-monotonicity} and
\eqref{eq:p071-stone-mollification-error} give
\begin{equation}
\label{eq:p071-stone-almost-monotone}
 w_{k+1}(z,x)
 \leq a_{k+1}(q(z),x)+\eta_{k+1}
 \leq a_k(q(z),x)+\eta_{k+1}
 \leq w_k(z,x)+\eta_{k+1}.
\end{equation}
Set
\begin{equation}
\label{eq:p071-stone-tail-correction}
                       c_k=\sum_{j\geq k+1}\eta_j,
             \qquad   g_k=w_k+c_k.
\end{equation}
Then \(c_k=\eta_{k+1}+c_{k+1}\), so
\eqref{eq:p071-stone-almost-monotone} implies \(g_{k+1}\leq g_k\).
Finally, \eqref{eq:p071-stone-envelope-convergence},
\eqref{eq:p071-stone-mollification-error}, and \(c_k\to0\) show
that \(g_k\downarrow v\) on \(C\times B\).  This proves the reverse
inclusion when \(d>0\).

If \(d=0\), the unique nonempty open subset of \(\mathbf R^0\) is a point.
Give \(Y\) a compatible metric and choose nested finite clopen partitions
\(\mathcal P_k\) whose mesh tends to zero, exactly as above.  Put
\begin{equation}
\label{eq:p071-stone-zero-dimensional-envelope}
        a_{k,P}=\max\{-k,\sup_{y\in P}\bar v(y)\},
        \qquad w_k(z)=a_{k,P_k(q(z))}.
\end{equation}
The functions \(w_k\) are locally constant, hence smooth psh.  Refinement
and the lower truncations show that they decrease.  Upper semicontinuity
and the shrinking mesh give \(w_k\downarrow v\), by the same argument as
for \eqref{eq:p071-stone-envelope-convergence}.  Thus the conclusion holds
without Euclidean mollification or a tail correction.

We have proved
\(\mathscr Q\subseteq\operatorname{Psh}_{X,1}\) in every dimension.
Together with \eqref{eq:p071-stone-minimality} and the tautological
inclusion
\(\operatorname{Psh}_{X,1}\subseteq\operatorname{Psh}_X\), this proves
\eqref{eq:p071-stone-collapse}--\eqref{eq:p071-stone-factorization}, as
well as the fixed-neighbourhood assertion
\eqref{eq:p071-stone-fixed-neighbourhood}.

For the displayed example, the Cantor cube \(\{0,1\}^{\Lambda}\) is compact,
Hausdorff, and zero-dimensional for every \(\Lambda\).  When \(\Lambda\) is
uncountable it is generally nonmetrizable, but every function produced by
the hierarchy locally factors through the metrizable quotient constructed
in \eqref{eq:p071-stone-common-quotient}; no metrizability of \(Z\) was
used.  This also makes clear why
the argument does not extend to a nontrivial affine structure in the
\(Z\)-direction or to general polyhedral face gluings.
